% Provenance (never printed): see sources/outline_sources.md and sources/book1/BRIEFS.md (g7/03).
\chapter{Strings}\label{ch:g7:strings}

Sami's favourite story is a few pages long, and he wonders which vowel it uses
most. Counting by hand, with a pencil and five tally marks, he loses his place
twice on the first page. For a \omterm{def:g4:programs-and-operating-system:program}{program}, a text is a row of letters, much like
a \omterm{def:g7:lists:list}{list} of numbers, and walking through it is something we already know how to
do.

\begin{recall}
Every \omterm{def:g7:strings:string}{character} has a number, its code, in a \omterm{def:g5:text-pictures-sound-as-numbers:code}{code table}, and text is stored as
those numbers (\cref{ch:g5:text-pictures-sound-as-numbers}). A \omterm{def:g7:lists:list}{list} has
indices from 0, a length given by \code{len}, and can be walked through with
\code{for} (\cref{ch:g7:lists}).
\end{recall}

\section{A string is a sequence of characters}

\begin{definition}[String and character]\label{def:g7:strings:string}
A \emph{string}\index{string} is a \omterm{def:g6:variables-and-input:value}{value} made of characters in order: a text,
written between quotes. A \emph{character}\index{character} is one letter,
digit, space or sign. Strings have the type \code{str}.
\end{definition}

A \omterm{def:g7:strings:string}{string} works like a \omterm{def:g7:lists:list}{list} of \omterm{def:g7:strings:string}{characters}: the same indices from 0, negative
indices from the end, and \code{len} for its length.

\Needspace*{8\baselineskip}
\omcode{code/grade-7/03-strings/listings/string_basics.py}{1}{5}{Indices and
length of a \omterm{def:g7:strings:string}{string}.}\label{lst:g7:strings:basics}

\Needspace*{8\baselineskip}
\omoutput{code/grade-7/03-strings/outputs/string_basics.txt}{\code{"!" * 3}
repeats a \omterm{def:g7:strings:string}{string}, as in \cref{ch:g6:variables-and-input}.}\label{out:g7:strings:basics}

\begin{omfigure}
\begin{tikzpicture}[x=1cm, y=1cm, font=\footnotesize]
  \foreach \c [count=\i from 0] in {P,y,t,h,o,n} {
    \draw[black!60, fill=ocYellow!20] (1.1*\i,0) rectangle ++(1.1,0.7);
    \node[font=\ttfamily] at (1.1*\i+0.55,0.35) {\c};
    \node[ocRed!80!black] at (1.1*\i+0.55,-0.3) {\i};
    \pgfmathtruncatemacro{\n}{\i-6}
    \node[ocBlue!80!black] at (1.1*\i+0.55,1.0) {\n};
  }
  \node[anchor=east, ocBlue!80!black] at (-0.2,1.0) {from the end};
  \node[anchor=east, ocRed!80!black] at (-0.2,-0.3) {index};
\end{tikzpicture}

{\small The \omterm{def:g7:strings:string}{string} \code{"Python"}: indices 0 to 5 from the start (below), $-1$
to $-6$ from the end (above).}
\end{omfigure}

Unlike a \omterm{def:g7:lists:list}{list}, a \omterm{def:g7:strings:string}{string} cannot be changed once it is made: Python refuses to
replace one of its \omterm{def:g7:strings:string}{characters}.

\Needspace*{5\baselineskip}
\omcode{code/grade-7/03-strings/listings/immutable.py}{1}{2}{Changing a
\omterm{def:g7:strings:string}{character}?}\label{lst:g7:strings:immutable}

\Needspace*{10\baselineskip}
\omoutput{code/grade-7/03-strings/outputs/immutable_run.txt}{A \omterm{def:g7:strings:string}{string} does not
support \omterm{def:g7:lists:list}{item} \omterm{def:g6:variables-and-input:assignment}{assignment}.}\label{out:g7:strings:immutable}

Instead, we build a new \omterm{def:g7:strings:string}{string} and give it the old name:

\Needspace*{6\baselineskip}
\omcode{code/grade-7/03-strings/listings/rebuild.py}{1}{3}{A new \omterm{def:g7:strings:string}{string} in place
of the old one.}\label{lst:g7:strings:rebuild}

\Needspace*{5\baselineskip}
\omoutput{code/grade-7/03-strings/outputs/rebuild.txt}{\code{name[1:]} is
``everything from \omterm{def:g7:lists:list}{index} 1 on''.}\label{out:g7:strings:rebuild}

\section{Slices}

\begin{definition}[Slice]\label{def:g7:strings:slice}
A \emph{slice}\index{slice} is a part of a \omterm{def:g7:strings:string}{string} (or of a \omterm{def:g7:lists:list}{list}) taken by
indices: \code{s[a:b]} holds the \omterm{def:g7:strings:string}{characters} from \omterm{def:g7:lists:list}{index} $a$ up to \omterm{def:g7:lists:list}{index} $b$, but
without $b$. Without $a$, the slice starts at the beginning; without $b$, it
goes to the end; \code{s[::-1]} walks backwards and gives the \omterm{def:g7:strings:string}{string} reversed.
\end{definition}

\Needspace*{9\baselineskip}
\omcode{code/grade-7/03-strings/listings/slices.py}{1}{6}{\omterm{def:g7:strings:slice}{Slices} of
\code{"keyboard"}.}\label{lst:g7:strings:slices}

\Needspace*{9\baselineskip}
\omoutput{code/grade-7/03-strings/outputs/slices.txt}{A \omterm{def:g7:strings:slice}{slice} never includes its
end \omterm{def:g7:lists:list}{index}.}\label{out:g7:strings:slices}

\begin{omfigure}
\begin{tikzpicture}[x=1cm, y=1cm, font=\footnotesize]
  \foreach \c [count=\i from 0] in {k,e,y,b,o,a,r,d} {
    \draw[black!60, fill=ocYellow!20] (1.0*\i,0) rectangle ++(1.0,0.7);
    \node[font=\ttfamily] at (\i+0.5,0.35) {\c};
  }
  \foreach \i in {0,...,8} \node[ocRed!80!black, font=\scriptsize] at (\i,-0.25) {\i};
  \draw[ocBlue, thick, decorate, decoration={brace, amplitude=4pt}] (0,0.85) -- (3,0.85) node[midway, above=4pt] {\code{word[0:3]}};
  \draw[ocMint!70!black, thick, decorate, decoration={brace, amplitude=4pt}] (3,0.85) -- (8,0.85) node[midway, above=4pt] {\code{word[3:]}};
  \draw[ocRed, thick, decorate, decoration={brace, amplitude=4pt, mirror}] (3,-0.5) -- (8,-0.5) node[midway, below=4pt] {\code{word[-5:]}};
\end{tikzpicture}

{\small Think of the indices as standing \emph{between} the \omterm{def:g7:strings:string}{characters}: the
\omterm{def:g7:strings:slice}{slice} \code{[a:b]} takes what lies between the marks $a$ and $b$.}
\end{omfigure}

\section{Walking through text}

\code{for ch in s:} visits the \omterm{def:g7:strings:string}{characters} one by one. Counting vowels is a walk
with a \omterm{def:g6:while-loops:counter}{counter}; building a new \omterm{def:g7:strings:string}{string} is a walk that adds \omterm{def:g7:strings:string}{characters} to a
\omterm{def:g7:strings:string}{string} that starts empty, \code{""}. The test \code{ch in "aeiouAEIOU"} is true
when the \omterm{def:g7:strings:string}{character} is one of those ten.

\omcode{code/grade-7/03-strings/listings/vowels.py}{1}{12}{Counting vowels and
removing spaces.}\label{lst:g7:strings:vowels}

\Needspace*{6\baselineskip}
\omoutput{code/grade-7/03-strings/outputs/vowels.txt}{18 vowels, and the text
without its spaces.}\label{out:g7:strings:vowels}

\section{Useful string methods}

\begin{definition}[String method]\label{def:g7:strings:string-method}
A \emph{string method}\index{string method} is a \omterm{def:g7:functions:function}{function} that belongs to
\omterm{def:g7:strings:string}{strings} and is called after a \omterm{def:g7:strings:string}{string} and a point, as in \code{s.upper()}. A
\omterm{def:g7:strings:string}{string} method never changes the \omterm{def:g7:strings:string}{string}: it returns a new \omterm{def:g6:variables-and-input:value}{value}.
\end{definition}

\Needspace*{11\baselineskip}
\omcode{code/grade-7/03-strings/listings/methods.py}{1}{8}{A few \omterm{def:g7:strings:string-method}{string
methods}.}\label{lst:g7:strings:methods}

\Needspace*{11\baselineskip}
\omoutput{code/grade-7/03-strings/outputs/methods.txt}{\code{strip} removes the
spaces at both ends; \code{find} gives the \omterm{def:g7:lists:list}{index} where \code{"World"} starts;
the last line shows that \code{s} is unchanged.}\label{out:g7:strings:methods}

Two methods turn a text into a \omterm{def:g7:lists:list}{list} of words and back: \code{split} cuts a
\omterm{def:g7:strings:string}{string} at each separator, and \code{join} glues the \omterm{def:g7:strings:string}{strings} of a \omterm{def:g7:lists:list}{list}, with a
separator between them. Without an argument, \code{split()} cuts at any run of
spaces and line breaks.

\Needspace*{9\baselineskip}
\omcode{code/grade-7/03-strings/listings/words.py}{1}{6}{Cutting a sentence into
words and gluing them back.}\label{lst:g7:strings:words}

\Needspace*{8\baselineskip}
\omoutput{code/grade-7/03-strings/outputs/words.txt}{A \omterm{def:g7:lists:list}{list} method,
\code{count}, counts the \omterm{def:g7:lists:list}{items} equal to
\code{"the"}.}\label{out:g7:strings:words}

\begin{omfigure}
\begin{tikzpicture}[x=1cm, y=1cm, font=\footnotesize\ttfamily]
  \node[draw=black!60, fill=ocYellow!20, minimum height=0.6cm] (s) at (3.2,0) {the fox and the hen};
  \foreach \w [count=\i from 0] in {the,fox,and,the,hen}
    \node[draw=black!60, fill=ocBlue!10, minimum width=1.0cm, minimum height=0.55cm] (w\i) at (1.3*\i+0.6,-1.4) {\w};
  \node[draw=black!60, fill=ocYellow!20, minimum height=0.6cm] (j) at (3.2,-2.8) {the-fox-and-the-hen};
  \draw[-{Stealth[length=2mm]}, thick] (s.south) -- node[right, font=\footnotesize\sffamily] {\code{split(" ")}: cut at each space} (3.2,-1.1);
  \draw[-{Stealth[length=2mm]}, thick] (3.2,-1.7) -- node[right, font=\footnotesize\sffamily] {\code{"-".join(...)}: glue with \code{-}} (j.north);
\end{tikzpicture}

{\small \code{split} cuts a \omterm{def:g7:strings:string}{string} into a \omterm{def:g7:lists:list}{list} of \omterm{def:g7:strings:string}{strings}; \code{join} glues a
\omterm{def:g7:lists:list}{list} of \omterm{def:g7:strings:string}{strings} into one \omterm{def:g7:strings:string}{string}.}
\end{omfigure}

\section{The code of a character}

\code{ord} gives the code of a \omterm{def:g7:strings:string}{character}, and \code{chr} gives the \omterm{def:g7:strings:string}{character}
of a code: Python uses the same codes as the table of
\cref{ch:g5:text-pictures-sound-as-numbers} for the letters, digits and signs
of English. Its table goes much further, and holds the \omterm{def:g7:strings:string}{characters} of every
alphabet: \code{chr(1046)} is a Cyrillic capital letter, Arabic and Hindi
\omterm{def:g7:strings:string}{characters} have codes in the thousands, and Chinese \omterm{def:g7:strings:string}{characters} in the tens of
thousands. The next book explains how
such codes are stored in \omterm{def:g4:memory-and-storage:byte}{bytes}.

\Needspace*{9\baselineskip}
\omcode{code/grade-7/03-strings/listings/codes.py}{1}{6}{Codes and
\omterm{def:g7:strings:string}{characters}.}\label{lst:g7:strings:codes}

\Needspace*{10\baselineskip}
\omoutput{code/grade-7/03-strings/outputs/codes.txt}{The small letters come 32
places after the capitals.}\label{out:g7:strings:codes}

Subtracting 32 turns a small \code{g} into a capital, but a capital \code{G}
becomes an apostrophe. \code{upper} knows which \omterm{def:g7:strings:string}{characters} have a capital and
leaves the others alone: it is always the better choice.

A palindrome reads the same in both directions. The \omterm{def:g7:functions:function}{function} below first puts
the text in small letters and removes its spaces, then compares it with its
reverse:

\Needspace*{10\baselineskip}
\omcode{code/grade-7/03-strings/listings/palindrome.py}{1}{7}{A palindrome
test.}\label{lst:g7:strings:palindrome}

\Needspace*{8\baselineskip}
\omoutput{code/grade-7/03-strings/outputs/palindrome.txt}{Capitals and spaces do
not count.}\label{out:g7:strings:palindrome}

Finally, a few counts on a short passage. A \omterm{def:g7:strings:string}{string} between triple quotes may
run over several lines; \code{ch.isalpha()} is true when the \omterm{def:g7:strings:string}{character} is a
letter.

\omcode{code/grade-7/03-strings/listings/passage_stats.py}{1}{15}{Counting
\omterm{def:g7:strings:string}{characters}, letters, words and vowels.}\label{lst:g7:strings:passage}

\Needspace*{12\baselineskip}
\omoutput{code/grade-7/03-strings/outputs/passage_stats.txt}{The line break
counts as a \omterm{def:g7:strings:string}{character}, too.}\label{out:g7:strings:passage}

\begin{omfigure}
\begin{tikzpicture}
  \begin{axis}[width=12cm, height=4.4cm, ybar, bar width=5pt, ymin=0,
      symbolic x coords={a,b,c,d,e,f,g,h,i,j,k,l,m,n,o,p,q,r,s,t,u,v,w,x,y,z}, xtick=data,
      tick label style={font=\scriptsize}, ylabel={count}, label style={font=\footnotesize},
      axis x line*=bottom, axis y line*=left]
    \addplot[fill=ocBlue!55, draw=none] table[x=letter, y=count] {code/grade-7/03-strings/figures/out/letters.dat};
  \end{axis}
\end{tikzpicture}

{\small How often each letter appears in the passage of
\cref{lst:g7:strings:passage}, capitals counted with small letters. A script
of the chapter counts them from the listing itself.}
\end{omfigure}

\section{Exercises}

\begin{exercise}[$\star$]\label{exo:g7:strings:1}
With \code{s = "computer"}, what are \code{s[1:4]}, \code{s[-1]} and
\code{s[::-1]}?
\end{exercise}

\begin{exercise}[$\star$]\label{exo:g7:strings:2}
What does \code{"a,b,c".split(",")} give? And
\code{" + ".join(["1", "2", "3"])}?
\end{exercise}

\begin{exercise}[$\star$]\label{exo:g7:strings:3}
What are \code{ord("C")} and \code{chr(100)}? Use the codes of
\cref{out:g7:strings:codes}.
\end{exercise}

\begin{exercise}[$\star$]\label{exo:g7:strings:4}
Write a \omterm{def:g4:programs-and-operating-system:program}{program} that counts the spaces of a sentence by walking through it,
and check the result with \code{count}.
\end{exercise}

\begin{exercise}[$\star\star$]\label{exo:g7:strings:5}
Tom writes \code{s[len(s)]} to get the last \omterm{def:g7:strings:string}{character} of a \omterm{def:g7:strings:string}{string}, and gets an
\code{IndexError}. Why? Give two correct ways.
\end{exercise}

\begin{exercise}[$\star\star$]\label{exo:g7:strings:6}
This palindrome test says that \code{"Kayak"} is not a palindrome. Why? Fix it.

\omcode{code/grade-7/03-strings/exercises/palindrome_bug.py}{1}{5}{A test that
gets capitals wrong.}\label{lst:g7:strings:palindrome-bug}
\end{exercise}

\begin{exercise}[$\star\star$]\label{exo:g7:strings:7}
Write a \omterm{def:g7:functions:function}{function} \code{initials(full_name)}: \code{"ada lovelace"} gives
\code{"A.L."}.
\end{exercise}

\begin{exercise}[$\star\star$]\label{exo:g7:strings:8}
Write a \omterm{def:g4:programs-and-operating-system:program}{program} that prints a text with every vowel replaced by \code{*}.
\end{exercise}

\begin{exercise}[$\star\star$]\label{exo:g7:strings:9}
How many different \omterm{def:g7:strings:slice}{slices} \code{s[i:i + 3]} of length 3 does a \omterm{def:g7:strings:string}{string} of length
10 have?
\end{exercise}

\begin{exercise}[$\star\star\star$]\label{exo:g7:strings:10}
Write a \omterm{def:g7:functions:function}{function} \code{longest_word(sentence)} that returns the longest word of a
sentence (the first one, if several are as long).
\end{exercise}

\begin{exercise}[$\star\star\star$]\label{exo:g7:strings:11}
Explain why \code{s[::-1][::-1]} is always \code{s}, and why
\code{s[:k] + s[k:]} is \code{s} for every \code{k} from 0 to \code{len(s)}.
\end{exercise}

\begin{exercise}[$\star\star\star$]\label{exo:g7:strings:12}
In a mirror, a word written in small letters appears reversed, with \code{b}
and \code{d} swapped, and \code{p} and \code{q} swapped. Write a \omterm{def:g7:functions:function}{function}
\code{mirror(word)} that returns what the mirror shows, using the two \omterm{def:g7:strings:string}{strings}
\code{"bdpq"} and \code{"dbqp"}, and find which of \code{bid}, \code{bud},
\code{pod}, \code{pq} and \code{box} look the same in the mirror.
\end{exercise}

\section{Problem: Palindrome Dates}

\begin{problem}[{Weekend problem --- dates that read both ways}]\label{pb:g7:strings:1}
Write a date with two digits for the day, two for the month and four for the
year, without separators: the 2nd of February 2020 is \code{02022020}. Some
dates read the same backwards: they are \emph{palindrome dates}.

\needspace{4\baselineskip}\medskip\noindent\textbf{Part I --- By hand.}
\begin{enumerate}
  \item Write 02-02-2020 and 12-02-2021 as eight digits, and check that both
    are palindromes.
  \item For a year from 2001 to 2099, written $2\,0\,c\,d$, the date must be
    $d\,c\,0\,2$ followed by the year. What does that say about the month of
    every palindrome date of this century?
  \item For the year 2092, which day would it have to be? Does that day exist?
\end{enumerate}

\needspace{4\baselineskip}\medskip\noindent\textbf{Part II --- Functions.}

\omcode{code/grade-7/03-strings/listings/dates.py}{1}{23}{Dates as \omterm{def:g7:strings:string}{strings} of
eight digits.}\label{lst:g7:strings:dates}

\begin{enumerate}[resume]
  \item Explain how \code{two} writes 7 as \code{"07"} and 12 as \code{"12"}.
  \item Why does \code{date_string} not call \code{two} for the year?
  \item Which chapter's rule does \code{days_in_month} use, and why does
    February need it?
  \item Run \code{dates.py}. Does it agree with question 1?
  \item Write a \omterm{def:g4:programs-and-operating-system:program}{program} that walks through every day of every year from 2001
    to 2099 and keeps the palindrome dates in a \omterm{def:g7:lists:list}{list}. (Use three nested
    \code{for} \omterm{def:g4:algorithms-repeat-and-decide:loop}{loops}.)
\end{enumerate}

\needspace{4\baselineskip}\medskip\noindent\textbf{Part III --- The list.}
\begin{enumerate}[resume]
  \item How many palindrome dates are there from 2001 to 2099? Which are the
    first and the last?
  \item Why is there exactly one palindrome date in some years, and none in
    others? Which years have one?
  \item Why is the last one so special?
  \item State the result.
\end{enumerate}
\end{problem}
